%% ============================================================
%%  SOLUCIÓN — Ejercicio 1.1: Macros propias
%%  Clase 2 — Después de diapositiva "Macros propias: nunca repitas código"
%%  Compilar: F5 (pdflatex)
%% ============================================================

\documentclass[12pt, a4paper]{article}
\usepackage[utf8]{inputenc}
\usepackage[T1]{fontenc}
\usepackage[spanish]{babel}
\usepackage{amsmath, amssymb}

% --- MACROS ECONOMÉTRICAS (en el preámbulo) ---
\newcommand{\E}{\mathbb{E}}
\newcommand{\Var}{\mathrm{Var}}
\newcommand{\Cov}{\mathrm{Cov}}
\newcommand{\plim}{\operatorname{plim}}
\newcommand{\R}{\mathbb{R}}
\newcommand{\deriv}[2]{\frac{\partial #1}{\partial #2}}
\newcommand{\abs}[1]{\left| #1 \right|}
\newcommand{\norm}[1]{\left\| #1 \right\|}

\title{Macros personalizadas en \LaTeX{}}
\author{Tu Nombre}
\date{\today}

\begin{document}
\maketitle

\section{Uso de macros en expresiones econométricas}

Con las macros definidas, escribimos expresiones de forma más limpia y consistente.

\subsection{Esperanza condicional}

\[
  \E[Y_i | X_i] = X_i'\beta
\]

Si queremos cambiar la notación de esperanza, solo editamos \verb|\E| una vez 
en el preámbulo, y todo el documento se actualiza automáticamente.

\subsection{Varianza del estimador}

\[
  \Var(\hat{\beta}) = \sigma^2(X'X)^{-1}
\]

\subsection{Covarianza entre variables}

\[
  \Cov(X_i, \varepsilon_i) = 0
\]

Este es el supuesto de exogeneidad estricta, crucial para insesgadez de OLS.

\subsection{Límite en probabilidad}

\[
  \plim \hat{\beta} = \beta
\]

Este es el resultado de consistencia del estimador OLS bajo supuestos estándar.

\subsection{Derivada parcial}

Usando la macro con argumentos:

\[
  \deriv{u(c)}{c} = c^{-\sigma}
\]

Esta es la utilidad marginal de consumo en un modelo estándar con preferencias 
de potencia.

\subsection{Valor absoluto y norma}

\[
  \abs{x} \quad \text{y} \quad \norm{\hat{\beta} - \beta} \leq C / \sqrt{N}
\]

La norma de la diferencia entre el estimador y el parámetro verdadero converge 
a cero a tasa $1/\sqrt{N}$.

\section*{Ventajas de las macros}

\begin{enumerate}
  \item \textbf{Consistencia}: todos los símbolos se ven iguales en todo el documento.
  \item \textbf{Mantenibilidad}: cambios globales requieren editar un único lugar.
  \item \textbf{Legibilidad}: el código es más limpio y semánticamente significativo.
  \item \textbf{Compatibilidad}: los comandos personalizados funcionan igual en 
        artículos, libros, presentaciones, etc.
\end{enumerate}

\end{document}
